\documentclass[11pt]{article}
\usepackage[margin=1in]{geometry}
\usepackage{listings}
\usepackage{xcolor}
\usepackage{booktabs}
\usepackage{hyperref}

\hypersetup{colorlinks=true, linkcolor=blue, urlcolor=blue}

\lstset{
  basicstyle=\ttfamily\small,
  keywordstyle=\color{blue},
  commentstyle=\color{gray},
  frame=single,
  breaklines=true,
  columns=fullflexible,
  numbers=left,
  numberstyle=\tiny\color{gray}
}

\title{The \texttt{arty.xdc} File, Explained Line by Line\\
\large \texttt{final\_8} project \,\textperiodcentered\, Digilent Arty A7-35T
\,\textperiodcentered\, includes the July 16--17 debug notes}
\author{Lucas Zhou}
\date{July 17, 2026}

\begin{document}
\maketitle
\tableofcontents
\newpage

%=====================================================================
\section{What Is an XDC File in the First Place?}

Your \texttt{top.v} describes \emph{logic}: ``there are inputs called
\texttt{A2, A1, B2, B1}, an output bus called \texttt{seg}, and here is the
truth table connecting them.'' Notice that nothing in the Verilog says
\emph{where} those signals physically live. The FPGA chip has 324 metal
balls on its underside, and the Arty board routes some of them to switches,
some to LEDs, some to the Pmod connectors on the edge.

The \textbf{XDC file} (\emph{Xilinx Design Constraints}) is the missing
link. It tells the place-and-route tool (\texttt{nextpnr-xilinx} in our
flow):

\begin{quote}
``The Verilog port named \texttt{A2} must come out on the physical pin
named \texttt{A10}, using 3.3\,V signaling.''
\end{quote}

That is essentially \emph{all} our XDC does --- eleven sentences of that
form, one per port. XDC files are written in the \textbf{Tcl} scripting
language (pronounced ``tickle''), which is why the syntax looks nothing
like Verilog. You never write algorithms in it here; every line is the
same one command with different arguments.

\textbf{Key mental model:} the XDC is a \emph{patch cable list}. Changing
it never changes your logic --- it only changes which physical pin each
logical wire pops out of. That is why every hardware bug in this project
was fixed in the XDC while your \texttt{top.v} stayed untouched.

%=====================================================================
\section{Anatomy of One Constraint Line}

Every functional line in the file is this single pattern:

\begin{lstlisting}[numbers=none]
set_property -dict { PACKAGE_PIN A10  IOSTANDARD LVCMOS33 } [get_ports { A2 }];
\end{lstlisting}

Reading it piece by piece:

\begin{description}
  \item[\texttt{set\_property}] The Tcl command. It means ``attach the
    following properties to the following object.'' Same command on every
    line.
  \item[\texttt{-dict \{ ... \}}] ``Here comes a \emph{dictionary} of
    property--value pairs.'' A dictionary is just an alternating list:
    \emph{name value name value}. Ours contains two pairs:
    \begin{itemize}
      \item \texttt{PACKAGE\_PIN A10} --- property \texttt{PACKAGE\_PIN}
        gets value \texttt{A10}. This is the physical ball on the chip
        package. Pin names like \texttt{A10}, \texttt{G13}, \texttt{E15}
        are grid coordinates on the 324-ball package (row letter + column
        number), \emph{not} related to your signal names. It is pure
        coincidence that a switch pin is called \texttt{A10} and your port
        is called \texttt{A2}.
      \item \texttt{IOSTANDARD LVCMOS33} --- property \texttt{IOSTANDARD}
        gets value \texttt{LVCMOS33}, meaning ``Low-Voltage CMOS,
        3.3\,V.'' FPGA pins can speak many electrical dialects (1.8\,V,
        2.5\,V, differential pairs, \dots). Everything on the Arty's
        switches and Pmod connectors is wired for 3.3\,V, so every line
        says \texttt{LVCMOS33}. Omitting it is an error in Vivado and bad
        practice everywhere.
    \end{itemize}
  \item[{\texttt{[get\_ports \{ A2 \}]}}] In Tcl, square brackets mean
    ``run this inner command first and substitute its result'' --- like
    $f(g(x))$ in math. \texttt{get\_ports} looks up a port \emph{by the
    name it has in the top-level Verilog module}. So this fetches your
    \texttt{input wire A2} and hands it to \texttt{set\_property}. The
    name must match the Verilog port exactly; for a bus you name one bit
    at a time: \texttt{seg[0]}, \texttt{seg[1]}, \dots
  \item[\texttt{;}] End of statement. (In Tcl a newline also ends a
    statement, so the semicolon is optional --- it is kept for clarity.)
  \item[\texttt{\#\# comment}] Anything after \texttt{\#} is a comment.
    Double \texttt{\#\#} is just a style convention, same as one.
\end{description}

So the whole line reads, in English: \emph{``Find the top-level port
named \texttt{A2}, solder it (logically) to package ball \texttt{A10}, and
drive it with 3.3\,V CMOS levels.''}

%=====================================================================
\section{The Current, Working XDC --- Every Line}

\begin{lstlisting}
## final_8: 4 switches in, seg[6:0] out to Pmod SSD straddling JA and JB
##
## The Pmod SSD has two 6-pin headers: J1 (AA-AD, GND, VCC) plugs into the
## TOP ROW of JA, and J2 (AE, AF, AG, CAT, GND, VCC) plugs into the TOP ROW
## of JB. So four of the signals live on JB pins 1-4, not JA's bottom row.
##
## The display glyphs are mounted 180 deg rotated as viewed on the Arty
## (decimal points sit top-left), so each module signal lights the segment
## diametrically opposite its name: AA = viewed bottom, AD = viewed top,
## AB = viewed lower-left, etc. The assignments below bake in both the
## JA/JB split and that rotation.
##
## top.v numbers segments 1-7 bottom-to-top (seg[0..6] = viewed D,C,E,G,B,F,A).
## If a segment still lights in the wrong position, permute these pin
## assignments rather than the top.v table.

## Switches (A2=SW3, A1=SW2, B2=SW1, B1=SW0)
set_property -dict { PACKAGE_PIN A10   IOSTANDARD LVCMOS33 } [get_ports { A2 }];
set_property -dict { PACKAGE_PIN C10   IOSTANDARD LVCMOS33 } [get_ports { A1 }];
set_property -dict { PACKAGE_PIN C11   IOSTANDARD LVCMOS33 } [get_ports { B2 }];
set_property -dict { PACKAGE_PIN A8    IOSTANDARD LVCMOS33 } [get_ports { B1 }];

## Seven-segment outputs (viewed position -> module signal -> FPGA pin)
set_property -dict { PACKAGE_PIN G13   IOSTANDARD LVCMOS33 } [get_ports { seg[0] }];  # viewed D (bottom)      -> AA -> JA1
set_property -dict { PACKAGE_PIN E16   IOSTANDARD LVCMOS33 } [get_ports { seg[1] }];  # viewed C (lower-right) -> AF -> JB2
set_property -dict { PACKAGE_PIN B11   IOSTANDARD LVCMOS33 } [get_ports { seg[2] }];  # viewed E (lower-left)  -> AB -> JA2
set_property -dict { PACKAGE_PIN D15   IOSTANDARD LVCMOS33 } [get_ports { seg[3] }];  # viewed G (middle)      -> AG -> JB3
set_property -dict { PACKAGE_PIN E15   IOSTANDARD LVCMOS33 } [get_ports { seg[4] }];  # viewed B (upper-right) -> AE -> JB1
set_property -dict { PACKAGE_PIN A11   IOSTANDARD LVCMOS33 } [get_ports { seg[5] }];  # viewed F (upper-left)  -> AC -> JA3
set_property -dict { PACKAGE_PIN D12   IOSTANDARD LVCMOS33 } [get_ports { seg[6] }];  # viewed A (top)         -> AD -> JA4

## Pmod SSD digit select (CAT on the module's J2 pin 4 = JB4, NOT JA10):
## picks which of the two digits is lit. top.v drives it constant.
set_property -dict { PACKAGE_PIN C15   IOSTANDARD LVCMOS33 } [get_ports { C }];
\end{lstlisting}

\subsection{Lines 1--16: the header comments}
Pure documentation --- the tool ignores them. They record the two hardware
facts discovered on July 16 (Section~\ref{sec:whyoldfailed}) so the next
reader does not rediscover them the hard way.

\subsection{Lines 18--22: the four switch inputs}
One line per input port. Each maps a Verilog input to one of the four
slide switches at the bottom edge of the board:

\begin{center}
\begin{tabular}{@{}llll@{}}
\toprule
Line & Verilog port & Package pin & Board part \\
\midrule
19 & \texttt{A2} (high bit of $A$) & \texttt{A10} & switch SW3 \\
20 & \texttt{A1} (low bit of $A$)  & \texttt{C10} & switch SW2 \\
21 & \texttt{B2} (high bit of $B$) & \texttt{C11} & switch SW1 \\
22 & \texttt{B1} (low bit of $B$)  & \texttt{A8}  & switch SW0 \\
\bottomrule
\end{tabular}
\end{center}

Where do pins like \texttt{A10} come from? From the board schematic ---
Digilent publishes a ``master XDC'' listing which package pin every
switch, LED and connector wire uses. Our project keeps a condensed copy in
\texttt{docs/hardware-reference.md}. These four lines have been correct
from day one and never changed.

\subsection{Lines 24--32: the seven segment outputs}
One line per bit of \texttt{output reg [6:0] seg}. Three name spaces
meet on each line, and the comments spell out all three:

\begin{enumerate}
  \item \textbf{Your numbering} (\texttt{seg[0]}..\texttt{seg[6]}): the
    problem's ``segment 1 to 7,'' which counts the segments
    \emph{bottom-to-top}: bottom bar, lower-right, lower-left, middle,
    upper-right, upper-left, top. This ordering was reverse-engineered
    from your own case table (e.g.\ your pattern for ``1'' sets exactly
    bits 1 and 4 $\Rightarrow$ those must be the two right-hand verticals).
  \item \textbf{Module signal names} (\texttt{AA}..\texttt{AG}): what the
    Pmod SSD's silkscreen calls its seven inputs.
  \item \textbf{FPGA package pins} (\texttt{G13}, \texttt{E16}, \dots):
    where the Arty routes each Pmod connector pin.
\end{enumerate}

Line by line:
\begin{center}
\begin{tabular}{@{}lllll@{}}
\toprule
Line & Port & Meaning in \texttt{top.v} & SSD signal & Pin (connector) \\
\midrule
25 & \texttt{seg[0]} & segment 1 = bottom bar    & \texttt{AA} & \texttt{G13} (JA pin 1) \\
26 & \texttt{seg[1]} & segment 2 = lower-right   & \texttt{AF} & \texttt{E16} (JB pin 2) \\
27 & \texttt{seg[2]} & segment 3 = lower-left    & \texttt{AB} & \texttt{B11} (JA pin 2) \\
28 & \texttt{seg[3]} & segment 4 = middle bar    & \texttt{AG} & \texttt{D15} (JB pin 3) \\
29 & \texttt{seg[4]} & segment 5 = upper-right   & \texttt{AE} & \texttt{E15} (JB pin 1) \\
30 & \texttt{seg[5]} & segment 6 = upper-left    & \texttt{AC} & \texttt{A11} (JA pin 3) \\
31 & \texttt{seg[6]} & segment 7 = top bar       & \texttt{AD} & \texttt{D12} (JA pin 4) \\
\bottomrule
\end{tabular}
\end{center}

Why the scrambled-looking middle column? Two physical quirks, both
explained in Section~\ref{sec:whyoldfailed}: the module spans \emph{two}
connectors (JA and JB), and its glyphs are mounted upside down, so signal
\texttt{AA} lights the \emph{bottom} segment as you view it, \texttt{AD}
the top, and so on.

\subsection{Lines 34--36: the digit select}
The Pmod SSD contains \emph{two} digits sharing the seven segment wires.
The \texttt{CAT} input chooses which digit is actually lit (in normal use
you would toggle it fast to show two different numbers; we hold it
constant from \texttt{assign C = 1'b1} in \texttt{top.v}, so one digit is
always on). It lives on JB pin 4 = package pin \texttt{C15}.

%=====================================================================
\section{The Old Versions, and Why Each Failed}
\label{sec:whyoldfailed}

The XDC went through three versions. The Verilog logic was correct the
whole time --- the July 13 simulation sweep of all 16 inputs matched the
case table exactly, and it still does. Every failure below is purely
about \emph{which physical pin each signal came out of}.

\subsection{Version 1 (July 13): straight-through guess}

\begin{lstlisting}
## Seven-segment outputs on Pmod JA
set_property -dict { PACKAGE_PIN G13  IOSTANDARD LVCMOS33 } [get_ports { seg[0] }];  # JA1 -> AA
set_property -dict { PACKAGE_PIN B11  IOSTANDARD LVCMOS33 } [get_ports { seg[1] }];  # JA2 -> AB
set_property -dict { PACKAGE_PIN A11  IOSTANDARD LVCMOS33 } [get_ports { seg[2] }];  # JA3 -> AC
set_property -dict { PACKAGE_PIN D12  IOSTANDARD LVCMOS33 } [get_ports { seg[3] }];  # JA4 -> AD
set_property -dict { PACKAGE_PIN D13  IOSTANDARD LVCMOS33 } [get_ports { seg[4] }];  # JA7 -> AE
set_property -dict { PACKAGE_PIN B18  IOSTANDARD LVCMOS33 } [get_ports { seg[5] }];  # JA8 -> AF
set_property -dict { PACKAGE_PIN A18  IOSTANDARD LVCMOS33 } [get_ports { seg[6] }];  # JA9 -> AG
set_property -dict { PACKAGE_PIN K16  IOSTANDARD LVCMOS33 } [get_ports { C }];       # JA10
\end{lstlisting}

This assumed (a) \texttt{seg[N]} simply means SSD signal number $N{+}1$
in alphabetical order \texttt{AA}$\to$\texttt{AG}, and (b) the whole
module plugs into connector JA alone, using both its rows.

\textbf{Symptom on hardware:} garbage shapes --- a ``7''-looking glyph,
disconnected fragments. \textbf{Wrong because} your \texttt{top.v}
numbers segments geometrically (bottom-to-top), not alphabetically, so
each lit segment landed in an unrelated position.

\subsection{Version 2 (July 17, first fix): permuted, still JA-only}

\begin{lstlisting}
set_property -dict { PACKAGE_PIN D12  IOSTANDARD LVCMOS33 } [get_ports { seg[0] }];  # JA4 -> AD
set_property -dict { PACKAGE_PIN A11  IOSTANDARD LVCMOS33 } [get_ports { seg[1] }];  # JA3 -> AC
set_property -dict { PACKAGE_PIN D13  IOSTANDARD LVCMOS33 } [get_ports { seg[2] }];  # JA7 -> AE
set_property -dict { PACKAGE_PIN A18  IOSTANDARD LVCMOS33 } [get_ports { seg[3] }];  # JA9 -> AG
set_property -dict { PACKAGE_PIN B11  IOSTANDARD LVCMOS33 } [get_ports { seg[4] }];  # JA2 -> AB
set_property -dict { PACKAGE_PIN B18  IOSTANDARD LVCMOS33 } [get_ports { seg[5] }];  # JA8 -> AF
set_property -dict { PACKAGE_PIN G13  IOSTANDARD LVCMOS33 } [get_ports { seg[6] }];  # JA1 -> AA
\end{lstlisting}

This version fixed the \emph{numbering} mismatch: decoding the case table
showed \texttt{seg[0..6]} must mean segments D, C, E, G, B, F, A (the
bottom-to-top order). Assuming the SSD's names matched the standard
seven-segment letters and everything sat on JA, this permutation should
have produced correct digits.

\textbf{Symptom on hardware:} still partial digits --- for input
\texttt{1011} ($2\times3$), only three segments of the expected six lit
up. \textbf{Wrong because} both physical assumptions from Version 1 were
still baked in, and both turned out to be false.

\subsection{What the photo revealed (the two real bugs)}

Photo \texttt{IMG\_8616} (July 17), input \texttt{1011}, expected ``6'':

\begin{enumerate}
  \item \textbf{The module straddles two connectors.} The Pmod SSD has
    two separate 6-pin headers: J1 (\texttt{AA AB AC AD GND VCC}) in the
    top row of JA, and J2 (\texttt{AE AF AG CAT GND VCC}) in the top row
    of \emph{JB}. Versions 1--2 drove \texttt{AE/AF/AG/CAT} on JA's
    \emph{bottom row} (pins \texttt{D13 B18 A18 K16}) --- pins the module
    does not even touch. Segments E, F, G and the digit select were
    floating: three of your seven wires went nowhere. That is exactly
    ``part of a number.''
  \item \textbf{The glyphs are mounted 180$^\circ$ rotated.} The decimal
    points sit at the \emph{top-left} of each digit (normally
    bottom-right) --- the giveaway that the LED package is upside down
    relative to its labels. The lit segments confirmed it: the FPGA was
    driving signals \texttt{AA}, \texttt{AC}, \texttt{AD} high, and the
    photo showed \emph{bottom}, \emph{upper-left}, \emph{top} lit ---
    each one diametrically opposite its textbook position. Uncorrected,
    every digit would render upside down (a 6 would read as a 9).
\end{enumerate}

A useful detail: the three segments that \emph{did} light were a
perfectly correct half of a ``6.'' That proved the logic and the
Version-2 permutation were right as far as JA reached, and isolated the
fault to the unwired JB half plus the rotation.

\subsection{Version 3 (current): why it works}

The final XDC applies two corrections on top of Version 2:

\begin{itemize}
  \item \textbf{JA/JB split:} \texttt{AE}$\to$\texttt{E15} (JB1),
    \texttt{AF}$\to$\texttt{E16} (JB2), \texttt{AG}$\to$\texttt{D15}
    (JB3), \texttt{CAT}$\to$\texttt{C15} (JB4). Now all eight module
    inputs are actually driven.
  \item \textbf{180$^\circ$ rotation:} swap each signal with its opposite
    (\texttt{AA}$\leftrightarrow$\texttt{AD},
    \texttt{AB}$\leftrightarrow$\texttt{AE},
    \texttt{AC}$\leftrightarrow$\texttt{AF}; \texttt{AG} is the middle
    bar and maps to itself). So ``viewed top'' is driven via \texttt{AD},
    ``viewed bottom'' via \texttt{AA}, etc.
\end{itemize}

Sanity check with input \texttt{1011} ($2\times3=6$): \texttt{top.v}
outputs \texttt{seg = 1101111}, i.e.\ every segment on except
\texttt{seg[4]} (upper-right) --- which is precisely the shape of a 6,
now rendered complete and upright.

%=====================================================================
\section{Debug Notes / Timeline (July 16--17)}

\begin{enumerate}
  \item \textbf{Simulation} (Icarus Verilog): all 16 input combinations
    produce \texttt{seg} patterns matching the case table --- consistent
    with the 2-bit $\times$ 2-bit multiplier (results 0,1,2,3,4,6,9).
    Logic verified; simulation cannot catch pin-mapping bugs.
  \item \textbf{Build + flash, XDC v1}: clean build, but the display
    showed scrambled shapes. Diagnosis: segment-numbering mismatch
    (geometric vs.\ alphabetical). Fix: permute pins (v2).
  \item \textbf{Build + flash, XDC v2}: still partial digits.
    \texttt{1011} gave three segments of a six.
  \item \textbf{Photo analysis} (\texttt{IMG\_8616}): revealed the
    JA+JB straddle (three segments + digit select never wired) and the
    180$^\circ$-rotated glyphs (decimal points top-left; each lit
    segment opposite its signal name).
  \item \textbf{Build + flash, XDC v3}: all seven segments and the digit
    select now driven, rotation corrected. Awaiting visual confirmation;
    if the wrong digit of the two is lit, flip \texttt{assign C = 1'b1}
    to \texttt{1'b0} in \texttt{top.v} and rebuild.
\end{enumerate}

\textbf{Lessons for next time:}
\begin{itemize}
  \item Simulation checks logic; only hardware (or a datasheet) checks
    pinout. They fail independently.
  \item When a peripheral misbehaves, first ask ``is every wire actually
    connected to something?'' Floating pins produce exactly the
    ``partial'' symptom seen here.
  \item Fix mapping problems in the XDC, not by rewriting a correct
    truth table. The Verilog never changed after July 13.
  \item A photo of a known input is a powerful debug tool: three lit
    segments were enough to deduce both remaining bugs.
\end{itemize}

\end{document}
